\chapter*{Abstract}
\label{cha:abstract-en}
\makeatletter

\addcontentsline{toc}{section}{\textbf{Abstract}}

\begin{center}
	\noindent \Large \textbf{\@titleEN}
\end{center}

\vspace*{8mm}
In the era of rapid development of artificial intelligence, voice and text assistant systems (such as Google Assistant, Amazon Alexa, or Apple Siri) have become an integral part of everyday life. Despite their popularity, commercial solutions of this type are characterized by significant limitations resulting from the closed nature of their ecosystems. Integrating new devices or adding custom functionalities typically requires utilizing dedicated SDKs, undergoing rigorous certification processes, and becoming entirely dependent on a third-party provider's cloud infrastructure. This results in a lack of interoperability between devices from different manufacturers and raises privacy concerns. Furthermore, existing systems do not offer a native mechanism for grouping diverse devices into a common, cohesive space controlled by a single command.

The response to the aforementioned problems is this engineering thesis, the main objective of which was to design and implement a distributed, modular, and fully open AI assistant system. The developed solution enables intuitive control over a heterogeneous set of end devices (including personal computers, smartphones, and smartwatches) using natural language commands. A fundamental design premise was the openness of the architecture, allowing for easy and dynamic extension of the system's capabilities to support new devices and domains without modifying its source code.

The architecture of the developed system is based on the client-server model. The central server, implemented in C\#, was designed according to the modular monolith paradigm with a rigorous application of the Domain-Driven Design (DDD) methodology and the CQRS (Command Query Responsibility Segregation) pattern. Such an architectural division allowed for the clear delineation of bounded contexts for individual system modules (e.g., managing connections, groups, plugins, and prompt processing). This approach mitigated the operational overhead characteristic of distributed microservice architectures while maintaining high cohesion and extensibility for future development. Bidirectional, low-latency communication between clients and the server is handled by the SignalR library, ensuring reliable, real-time data transport.

A crucial element ensuring the system's extensibility is the innovative plugin architecture on the client side. Device-specific functionalities are delivered as dynamically loaded .dll libraries. Upon startup, the client application loads available modules into memory and registers their metadata with the central server, leaving the execution code exclusively on the user's device. Consequently, the server acts as a stateless orchestrator, lacking knowledge of the implementation details of individual operations.

The central logical entity of the designed solution is device groups. This mechanism allows the user to connect multiple physical machines based on a generated access code, creating a unified operational space. When a user issues a command (prompt), the server aggregates the metadata of all active plugins within the specified group and constructs an extended context for the Large Language Model (LLM). Integration with AI models was achieved using the Microsoft Semantic Kernel library, which serves as an abstraction layer, decoupling the business logic from any specific LLM provider. By utilizing the function calling mechanism, the model analyzes the user's intent and returns the optimal set of functions to execute. Following the user's authorization of the operation, the server distributes commands to the target clients, where they are executed natively.

The developed system constitutes an open smart assistant ecosystem. The designed architecture effectively addresses the problems of interoperability and extensibility, enabling the collaboration of numerous diverse devices. Utilizing modern software engineering patterns in conjunction with the power of Large Language Models has enabled the creation of a tool that serves as a scalable alternative to commercial solutions, ensuring the user's independence and full control over their digital environment.
